\section{Gyroscopic virtual mechanical system}\label{sec:model}

\subsection{Nominal interconnection}
The host is a normalised undamped oscillator with displacement $q$; an impulse
sets $(q,\dot q)(0)=(0,v_0)$.
The virtual oscillator, with state $(z,\dot z)$, lives in software.
The nominal interconnection is
\begin{equation}\label{eq:nominal}
  \ddot q - n\dot z + q = 0,\qquad \ddot z + n\dot q + z = 0,
\end{equation}
with gyroscopic coupling $n>0$.
The force $n\dot z$ is applied to the host by an actuator; the input $-n\dot q$ to
the virtual oscillator is computed from the measured host velocity.
With $x=(q,\dot q,z,\dot z)\in\RR^4$, \eqref{eq:nominal} reads $\dot x=A_cx$, and
the total energy $H(x)=\tfrac12|x|^2$ satisfies $\dot H=n\dot q\dot z-n\dot z\dot
q=0$.
The actuator injects power $n\dot q\dot z$ into the host and the virtual channel
withdraws the same amount, so the interconnection is lossless and
power-preserving~\citep{vanderschaft2000l2}.
The subsystem energies are $H_q=\tfrac12(q^2+\dot q^2)$ and
$H_z=\tfrac12(z^2+\dot z^2)$, and $\Hr:=H(x(0))=v_0^2/2$.

\begin{remark}[Sign convention]
\citet{juchem2026shaping} writes the coupling as $\ddot q+n\dot z+q=0$,
$\ddot z-n\dot q+z=0$; the change $z\mapsto-z$ maps one convention to the other
and leaves every statement about $(q,\dot q)$ unchanged.
We follow~\citet{juchem2027restoring}.
\end{remark}

The characteristic polynomial of $A_c$ is $(s^2+1)^2+n^2s^2$, with roots
$\pm\ii\Op$ and $\pm\ii\Om$, where
\begin{equation}\label{eq:modal-freq}
  \Opm=\tfrac12\big(\sqrt{n^2+4}\pm n\big),\quad \Op\Om=1,\quad \Op-\Om=n,
\end{equation}
and $S:=\Op+\Om=\sqrt{n^2+4}$.

\subsection{Exact modal--polar coordinates}
The following coordinates are the backbone of the paper.
They are defined on all of $\RR^4$, not only along nominal trajectories.

\begin{lemma}[Modal--polar coordinates]\label{lem:modal}
Define the linear map $T:\RR^4\to\RR^4$, $x\mapsto(X_+,Y_+,X_-,Y_-)$, by
\begin{equation}\label{eq:T}
\begin{aligned}
  X_+&=\tfrac{1}{S}(\dot q+\Om z), & Y_+&=\tfrac{1}{S}(\Om q-\dot z),\\
  X_-&=\tfrac{1}{S}(\dot q-\Op z), & Y_-&=\tfrac{1}{S}(\Op q+\dot z),
\end{aligned}
\end{equation}
and $\zeta_\pm:=X_\pm+\ii Y_\pm=:A_\pm e^{\ii\theta_\pm}$.
Then:
\begin{enumerate}
\item $T$ is invertible, with $q=Y_++Y_-$, $\dot q=\Op X_++\Om X_-$,
  $z=X_+-X_-$, $\dot z=-\Op Y_++\Om Y_-$.
\item Along \eqref{eq:nominal}, $\dot\zeta_\pm=\ii\Opm\zeta_\pm$, i.e.\
  $\dot A_\pm=0$ and $\dot\theta_\pm=\Opm$.
\item $H=\tfrac S2\big(\Op A_+^2+\Om A_-^2\big)$.
\item $m:=q\dot q+z\dot z = S A_+A_-\sin(\theta_++\theta_-)$.
\item $g:=\dot q\dot z=-\tfrac12\Op^2A_+^2\sin2\theta_+
  +\tfrac12\Om^2A_-^2\sin2\theta_- + A_+A_-\sin(\theta_--\theta_+)$.
\end{enumerate}
\end{lemma}
\begin{proof}
Items 1 and 2 follow by direct substitution, using $\Op\Om=1$ and $S=\Op+\Om$.
For item 3, item 1 gives $q^2+\dot z^2=(1+\Op^2)Y_+^2+(1+\Om^2)Y_-^2+2(1-\Op\Om)Y_+Y_-$
and $1+\Op^2=\Op S$, and similarly for $\dot q^2+z^2$.
Items 4 and 5 follow by substituting item 1 and applying product-to-sum
identities.
\end{proof}

Item 1 with $\theta_\pm=\Opm t+\phi_\pm$ recovers the nominal modal
representation of~\citet{juchem2027restoring}.
The impulse response has $A_+=A_-=v_0/S$ and $\theta_\pm(0)=0$.
The three \emph{error coordinates} used throughout are
\begin{equation}\label{eq:errors}
\begin{gathered}
  \eH:=\frac{H}{\Hr}-1,\qquad \eA:=\ln\frac{A_+}{A_-},\\
  \psi:=\sigma\theta_+-\tau\theta_-\ \ (\mathrm{mod}\ 2\pi),
\end{gathered}
\end{equation}
where $\eA$ and $\psi$ are defined on $\mathcal D_0:=\{x:A_+A_->0\}$ and
$(\tau,\sigma)$ is the design pair of \cref{sec:tuning}.
By item 3, $(H,\eA)$ determine $(A_+,A_-)$; the energy share of the $+$ mode is
\begin{equation}\label{eq:share}
  a:=\frac{S\Op A_+^2}{2H}=\frac{\Op e^{2\eA}}{\Op e^{2\eA}+\Om}\in(0,1).
\end{equation}

\begin{lemma}[Perturbed modal dynamics]\label{lem:perturbed}
Let $\ddot q=-q+n\dot z+\Delta_1$ and $\ddot z=-z-n\dot q+\Delta_2$ for arbitrary
(possibly state-, time- or history-dependent) signals $\Delta_1,\Delta_2$.
Then
\begin{equation}\label{eq:pert-modal}
  \dot\zeta_\pm=\ii\Opm\zeta_\pm+\tfrac{1}{S}\big(\Delta_1\mp\ii\Delta_2\big),
\end{equation}
and consequently
\begin{align}
  \dot H &= \dot q\,\Delta_1+\dot z\,\Delta_2, \label{eq:Hdot}\\
  \tfrac{\dd}{\dd t}\tfrac{A_\pm^2}{2} &= \tfrac1S\big(X_\pm\Delta_1\mp Y_\pm\Delta_2\big),
  \label{eq:Adot}\\
  \dot\theta_\pm &= \Opm\mp\tfrac{1}{SA_\pm^2}\big(X_\pm\Delta_2\pm Y_\pm\Delta_1\big).
  \label{eq:thdot}
\end{align}
\end{lemma}
\begin{proof}
Differentiate \eqref{eq:T}: $\dot X_\pm$ gains $\Delta_1/S$, $\dot Y_+$ gains
$-\Delta_2/S$ and $\dot Y_-$ gains $+\Delta_2/S$; the nominal part is item 2 of
\cref{lem:modal}.
Then \eqref{eq:Adot}--\eqref{eq:thdot} follow from
$\tfrac{\dd}{\dd t}|\zeta|^2=2\Real(\bar\zeta\dot\zeta)$ and
$\dot\theta=\Imag(\bar\zeta\dot\zeta)/|\zeta|^2$, and \eqref{eq:Hdot} from
item 3 of \cref{lem:modal} (or directly from $H=\tfrac12|x|^2$).
\end{proof}

\Cref{lem:perturbed} is exact and needs no averaging.
Every imperfection in \cref{sec:mismatch} and every control action in
\cref{sec:control} enters through $(\Delta_1,\Delta_2)$, and
\eqref{eq:Hdot}--\eqref{eq:thdot} show directly how it affects energy, amplitudes
and phases.
For example, a common coupling change $c$ ($\Delta_1=c\dot z$,
$\Delta_2=-c\dot q$) gives $\dot H=0$ and
\begin{equation}\label{eq:common-amp}
  \dot A_+=\tfrac{c}{S}\Om A_-\sin\varphi,\qquad
  \dot A_-=-\tfrac{c}{S}\Op A_+\sin\varphi,
\end{equation}
with $\varphi:=\theta_++\theta_-$.

\begin{remark}[Corrections to the conference version]\label{rem:ecc-fix}
The second identity in \eqref{eq:common-amp} carries a minus sign; it is
required by $\dot H=0$ (\cref{lem:modal}, item 3).
In~\citet[eq.~(48)]{juchem2027restoring} both amplitude rates were printed with
the same sign, although the derived quantity $B_A$ in their eq.~(50) is correct.
Likewise, $A_+=A_-$ with $H=\Hr$ gives $A_\pm=\sqrt{2\Hr}/S$, not
$\sqrt{2\Hr/S}$ as printed in~\citet[eq.~(52)]{juchem2027restoring}.
\end{remark}
